\begin{resumo}
Este trabalho busca propor uma análise comparativa entre duas arquiteturas de processamento
para sistemas de telemetria de hidrômetros analógicos: a computação em borda
(\textit{Edge Computing}), com inferência de \gls{ocr} realizada localmente no microcontrolador
\gls{esp32cam} por meio da biblioteca TensorFlow Lite for Microcontrollers, e a computação
em nuvem (\textit{Cloud Computing}), com transmissão da imagem capturada para
processamento remoto pelo motor Tesseract em conjunto com a biblioteca OpenCV.
A motivação reside na problemática do desperdício hídrico no Brasil e no elevado custo de substituição dos medidores
convencionais por modelos inteligentes nativos, o que torna o \textit{\gls{retrofitting}} de
hidrômetros analógicos uma alternativa economicamente viável. A transmissão dos dados
entre os nós de borda e o \textit{gateway} central é realizada via rede \gls{lorawan},
utilizando o Spreading Factor SF8, configuração que assegura conformidade com o limite
regulatório de \textit{dwell time} de 400~ms estabelecido pelo Ato n\textordmasculine~14.448
da ANATEL, ao mesmo tempo em que maximiza o alcance dentro desse limite. A viabilidade
do envio de imagens fragmentadas via \gls{lorawan} foi demonstrada analiticamente, com
uma imagem de 160$\times$120 pixels em escala de cinza, no formato JPEG com
aproximadamente 3,10~KB, podendo ser transmitida em 38 pacotes em cerca de 13,3~s
de rádio ativo, consumindo apenas 1,5\% do orçamento diário no cenário regulatório
mais restritivo (EU868 com \textit{duty cycle} de 1\%). Os dois fluxos de processamento
serão avaliados segundo quatro métricas quantitativas: acurácia de leitura \gls{ocr} (\%),
latência de ponta a ponta (ms), consumo energético por leitura (mAh) e volume de dados
transmitidos (bytes/leitura). Espera-se que os resultados orientem a escolha arquitetural
mais adequada para implantações de \textit{Smart Metering} em larga escala, considerando
o \textit{trade-off} entre eficiência energética, precisão de reconhecimento e complexidade
de manutenção do sistema.
\vspace{\onelineskip}

\noindent
\textbf{Palavras-chave}: Telemetria. ESP32. LoRaWan. Regressão. OCR.
\end{resumo}
